% GENERATED FILE. Edit canonical lessons, then run npm run generate:books.
% canonical-source-hash: fnv1a64:9822df67521c44ab
% canonical-lessons: ES-C362-tierra, ES-C362-mar, ES-C362-aire, ES-C362-fuego, ES-C362-luna

\chapter{Earth, Sea And Sky}
\label{ch:la-tierra-y-el-mar}
\hlchaptermodality{362}

\begin{chapteropening}
\textbf{By the end of this chapter:} \emph{I can say la tierra, el mar, el aire, el fuego and la luna, and say which of them a place is near.}

Complete ES-C362-luna: place a thing against earth, sea, air and fire, and name what is overhead at night.
\end{chapteropening}

\section[la tierra]{la tierra — the earth, and a mistake the Romans made about it}

\label{lesson:ES-C362-tierra}



\begin{quote}
Say \emph{the salt}, then \emph{the oil}. (\emph{La sal}; \emph{el aceite}.)
\end{quote}

\subsection*{What to know first}
\begin{itemize}
\raggedright
  \item el campo — the worked land, which is one piece of \emph{la tierra}.
  \item el mundo — the world entire, the wider word this one sits inside.
\end{itemize}

\begin{sounds}
\begin{itemize}
\raggedright
  \item \emph{TIE-rra}, two beats with the weight on the first. The \emph{rr} is rolled, held a moment longer than a single tap. $\to$ pronunciation reference
\end{itemize}
\end{sounds}

\begin{cousinweb}
\textbf{la tierra} — \textquotedblleft{}the earth, the ground.\textquotedblright{}

Latin \textbf{terra} meant the dry part of the world, the half that is not water. English borrowed it hard: \textbf{terrain}, \textbf{terrace}, \textbf{subterranean}, \textbf{terrestrial}, and \textbf{terra cotta}, which is \textquotedblleft{}cooked earth\textquotedblright{} and describes what a flowerpot is with no room left over. The \textbf{Mediterranean} is the sea \emph{in the middle of the land}, and to \textbf{inter} someone is to put them into the earth.

The best of the family is small and loud. A \textbf{terrier} is the dog bred to go into the ground after a badger — an animal named for the place it works.

Now the trap. Romans themselves thought \textbf{territōrium} came from \textbf{terrēre}, \textquotedblleft{}to frighten\textquotedblright{}: territory as the ground you could scare intruders off. They had their own etymology wrong. \emph{Territōrium} is \textbf{terra} with an ending meaning \textquotedblleft{}the land around a town\textquotedblright{}, and there is nothing frightening in it.

\emph{Terrēre} is real, though, and it left descendants of its own: \textbf{terror}, \textbf{terrible}, \textbf{terrific}. That is the frightening family, and \emph{tierra} is no part of it. The two roots even collide inside Spanish, where \emph{aterrar} is two different verbs wearing one spelling — one that terrifies, one that brings a thing down to the ground. The word you will hear at an airport, \emph{aterrizar}, \textquotedblleft{}to land\textquotedblright{}, belongs to the earth one.
\end{cousinweb}

\begin{grammarlens}[title={What you've built}]
A place-word wide enough to hold everything that follows: \emph{la tierra}, with \emph{el campo} inside it and \emph{el mundo} around it.
\end{grammarlens}

\subsection*{Your turn}
Say these aloud:

\begin{itemize}
\raggedright
  \item \emph{la tierra}
  \item \emph{la tierra y el campo}
  \item \emph{el mundo, la tierra, el campo}
\end{itemize}

\begin{tcolorbox}[breakable,colback=teal!4,colframe=teal!35!black,title={Before you move on}]
What did Roman scholars wrongly think territory came from? (\emph{Terrēre}, to frighten.) And which of the two roots does \emph{terror} belong to? (\emph{Terrēre}, not \emph{terra}.)
\end{tcolorbox}
\section[el mar]{el mar — the sea, and the mermaid who is a cousin not a copy}

\label{lesson:ES-C362-mar}



\begin{quote}
Say \emph{the oil}, then \emph{the earth}. (\emph{El aceite}; \emph{la tierra}.)
\end{quote}

\subsection*{What to know first}
\begin{itemize}
\raggedright
  \item el agua — the substance the sea is made of.
  \item la orilla — the edge where the sea has to stop.
\end{itemize}

\begin{sounds}
\begin{itemize}
\raggedright
  \item \emph{MAR}, one beat only. The \emph{r} at the end is a single tap of the tongue, not a roll. $\to$ pronunciation reference
\end{itemize}
\end{sounds}

\begin{cousinweb}
\textbf{el mar} — \textquotedblleft{}the sea.\textquotedblright{}

Latin \textbf{mare} came from an ancient root that meant sea long before Latin existed. Out of Latin, English borrowed a whole shelf of formal words: \textbf{marine}, \textbf{maritime}, \textbf{marina}, \textbf{submarine}. \textbf{Ultramarine} is the oddest of them — the blue \textquotedblleft{}from beyond the sea\textquotedblright{}, because the lapis it was ground from arrived by ship from Afghanistan.

Those are borrowings. Now the other route. English inherited that very same ancient root down its own line, where it surfaced as Old English \textbf{mere}, a lake or a sea. The word is still hiding inside \textbf{mermaid}: a \emph{mere}-maid, a maiden of the sea. \textbf{Marsh} and \textbf{moor} came down the same way, both of them waterlogged ground.

So \emph{mar} and the \emph{mer-} of \emph{mermaid} are not one word taken from the other. They are cousins, each walking down from a shared ancestor by a different road — one road through Rome, the other through the north.

One more, and it is a case of English mishearing itself. Latin called the shrub \emph{rōs marīnus}, standardly explained as \textquotedblleft{}dew of the sea\textquotedblright{}, for a plant that grows in salt wind. English bent that strange phrase into two familiar words and got \textbf{rosemary}, a rose and a Mary. It is not a rose, and it owes nothing to anyone named Mary. Spanish came down from the same \emph{rōs marīnus} and never pulled it toward either — the reshaping happened on the English side alone.
\end{cousinweb}

\begin{grammarlens}[title={What you've built}]
A word for the water with no edge to it, ready to set against the land: \emph{el mar}, \emph{el agua}, \emph{la orilla}.
\end{grammarlens}

\subsection*{Your turn}
Say these aloud:

\begin{itemize}
\raggedright
  \item \emph{el mar}
  \item \emph{el agua del mar}
  \item \emph{la orilla del mar}
\end{itemize}

\begin{tcolorbox}[breakable,colback=teal!4,colframe=teal!35!black,title={Before you move on}]
Is English \emph{mermaid} borrowed from Latin \emph{mare}, or inherited? (Inherited — the \emph{mer-} is Old English \emph{mere}, a cousin of \emph{mare}.) And what does \emph{ultramarine} mean word for word? (Beyond the sea, where the blue stone was shipped from.)
\end{tcolorbox}
\section[el aire]{el aire — the air that named a disease wrongly}

\label{lesson:ES-C362-aire}



\begin{quote}
Say \emph{the earth}, then \emph{the sea}. (\emph{La tierra}; \emph{el mar}.)
\end{quote}

\subsection*{What to know first}
\begin{itemize}
\raggedright
  \item el viento — air once it has a direction and a speed.
  \item el cielo — the sky that the air fills.
\end{itemize}

\begin{sounds}
\begin{itemize}
\raggedright
  \item \emph{AI-re}, two beats with the weight on the first. The \emph{ai} is one glide, the sound of English \emph{eye}. $\to$ pronunciation reference
\end{itemize}
\end{sounds}

\begin{cousinweb}
\textbf{el aire} — \textquotedblleft{}the air.\textquotedblright{}

Two journeys made this word. Greek had \textbf{aēr}, the lower air you breathe, as against the bright upper air of the gods. Latin borrowed it whole. Spanish then inherited it from Latin. Greek to Rome to Madrid, no stop in between.

English took the Greek shape for its technical vocabulary: \textbf{aerial}, \textbf{aerobic}, \textbf{aerosol}, \textbf{aeroplane}. All of them are breathing the same syllable.

The finest cousin is a mistake nobody has bothered to undo. Italian \emph{mala aria} meant \textquotedblleft{}bad air\textquotedblright{}, and before anyone knew what a mosquito did, the fever that haunted marshes was blamed on the vapour rising off them. English borrowed the phrase in the eighteenth century and welded it into one word, \textbf{malaria}. The theory is dead; the word carries it around regardless.

Here is the pleasing part. Spanish mostly calls that same disease by a name built on Latin \emph{palūs}, \textquotedblleft{}a swamp\textquotedblright{}. Spanish blamed the marsh, English blamed the air, and each language fossilised a different half of one abandoned idea.

Last, a small one. Italian \textbf{aria}, the piece a singer stands up to deliver, is this word yet again: \emph{air} in the sense of a tune. English used \emph{air} that way for centuries, which is why old songbooks are full of them.
\end{cousinweb}

\begin{grammarlens}[title={What you've built}]
A third thing to place a body in, after ground and water: \emph{el aire}, with \emph{el viento} moving through it and \emph{el cielo} over it.
\end{grammarlens}

\subsection*{Your turn}
Say these aloud:

\begin{itemize}
\raggedright
  \item \emph{el aire}
  \item \emph{el aire y el viento}
  \item \emph{el cielo, el aire, el viento}
\end{itemize}

\begin{tcolorbox}[breakable,colback=teal!4,colframe=teal!35!black,title={Before you move on}]
What did \emph{malaria} mean before it meant a disease? (Bad air, in Italian.) And which language did Latin take \emph{aer} from? (Greek.)
\end{tcolorbox}
\section[el fuego]{el fuego — the hearth that burned down the word for fire}

\label{lesson:ES-C362-fuego}



\begin{quote}
Say \emph{the sea}, then \emph{the air}. (\emph{El mar}; \emph{el aire}.)
\end{quote}

\subsection*{What to know first}
\begin{itemize}
\raggedright
  \item el calor — the heat a fire throws off.
  \item el horno — a fire with walls built around it.
\end{itemize}

\begin{sounds}
\begin{itemize}
\raggedright
  \item \emph{FUE-go}, two beats with the weight on the first. The \emph{g} is hard, the \emph{g} of English \emph{go}. $\to$ pronunciation reference
\end{itemize}
\end{sounds}

\begin{cousinweb}
\textbf{el fuego} — \textquotedblleft{}the fire.\textquotedblright{}

Latin had a perfectly serviceable word for fire: \textbf{ignis}, which survives in English \textbf{ignite} and \textbf{igneous}. Spanish threw it away.

What Spanish kept instead was \textbf{focus}, and in Latin \emph{focus} did not mean fire at all. It meant the hearth — the built place in a house where the fire was kept. The same swap ran right across Romance: French \emph{feu}, Italian \emph{fuoco}, Spanish \emph{fuego}, each of them a hearth that took over the work of the flame.

English then borrowed the hearth back, three separate times, and the image survives in each of them:

\begin{itemize}
\raggedright
  \item \textbf{focus} was Kepler's word, in 1604, for the burning point of a lens or a curved mirror. Put a hearth at the spot where the heat gathers and you have the picture exactly. \textbf{Focal} followed on from it.
  \item \textbf{curfew} is Old French \emph{cuevrefeu}, \textquotedblleft{}cover the fire\textquotedblright{} — the evening bell that told a town to bank its embers before bed.
  \item \textbf{fuel} is Old French \emph{fouaille}, from Late Latin \emph{focālia}, the stuff you feed a hearth with.
\end{itemize}

Spanish went further than any of that. It let the hearth become the word for home itself, so the fireplace ended up naming the house around it. English did something similar when it made \emph{hearth} and \emph{home} neighbours — two languages agreeing on where a house really begins.
\end{cousinweb}

\begin{grammarlens}[title={What you've built}]
The fourth element and the hottest: \emph{el fuego}, the \emph{el calor} it gives off, and \emph{el horno} that keeps it in one place.
\end{grammarlens}

\subsection*{Your turn}
Say these aloud:

\begin{itemize}
\raggedright
  \item \emph{el fuego}
  \item \emph{el calor del fuego}
  \item \emph{el fuego en el horno}
\end{itemize}

\begin{tcolorbox}[breakable,colback=teal!4,colframe=teal!35!black,title={Before you move on}]
Which Latin word is \emph{fuego} from, and what did it mean? (\emph{Focus}, the hearth.) And what does \emph{curfew} literally tell a town to do? (Cover the fire.)
\end{tcolorbox}
\section[la luna]{la luna — the shiny one, and a root you already own}

\label{lesson:ES-C362-luna}



\begin{quote}
Say \emph{the air}, then \emph{the fire}. (\emph{El aire}; \emph{el fuego}.)
\end{quote}

\subsection*{What to know first}
\begin{itemize}
\raggedright
  \item la luz — the same ancient root as this word, met earlier under another shape.
  \item la noche — when the moon does its work.
\end{itemize}

\begin{sounds}
\begin{itemize}
\raggedright
  \item \emph{LU-na}, two beats with the weight on the first. The \emph{u} is a pure \emph{oo}, never the \emph{yoo} of English \emph{lunar}. $\to$ pronunciation reference
\end{itemize}
\end{sounds}

\begin{cousinweb}
\textbf{la luna} — \textquotedblleft{}the moon.\textquotedblright{}

You already own this root. When \emph{la luz} was taught, the source was Latin \textbf{lūx}, light. \emph{Luna} is that same root wearing a different ending, and the moon turns out to be named, quite literally, \textbf{the shiny one}.

That is not a guess. Latin \emph{lūna} is worn down from an earlier \textbf{louksna}, built on the ancient root \textbf{lewk-}, \textquotedblleft{}to shine\textquotedblright{} — and archaic Latin actually preserves the older shape, spelled \textbf{Losna}, scratched onto a bronze mirror from Praeneste. The relatives elsewhere agree: Old Church Slavonic \emph{luna}, Old Prussian \emph{lauxnos} meaning \textquotedblleft{}stars\textquotedblright{}, Avestan \emph{raoxshna}.

Out of that one root, English got three sets of words through three different doors:

\begin{itemize}
\raggedright
  \item \textbf{light} walked in. It is the inherited Germanic cousin, borrowed from nobody.
  \item \textbf{lucid}, \textbf{luminous}, \textbf{illuminate}, \textbf{translucent} were carried in from Latin by scholars.
  \item \textbf{leukemia} arrived through Greek, where the root came out as \emph{leukós}, \textquotedblleft{}white\textquotedblright{} — a disease named for pale blood.
\end{itemize}

\textbf{Lunatic} belongs here too, from the long-held belief that the moon unsettled the mind.

Then the payoff lands on a Monday. Spanish \emph{lunes} is Latin \emph{diēs Lūnae}, \textquotedblleft{}day of the Moon\textquotedblright{} — itself a translation, piece for piece, of a Greek phrase. English \textbf{Monday} is the same translation, made independently into Germanic. Two languages, one Greek habit, one moon.
\end{cousinweb}

\begin{grammarlens}[title={What you've built}]
The thing overhead at night, tied back to a word you already met by day: \emph{la luna} and \emph{la luz}, one root in two shapes.
\end{grammarlens}

\subsection*{Your turn}
Say these aloud:

\begin{itemize}
\raggedright
  \item \emph{la luna}
  \item \emph{la luz de la luna}
  \item \emph{la luna en la noche}
\end{itemize}

\begin{tcolorbox}[breakable,colback=teal!4,colframe=teal!35!black,title={Before you move on}]
What does the name \emph{luna} mean word for word? (The shiny one.) And which archaic Latin form confirms the older shape? (\emph{Losna}, scratched on the Praeneste mirror.)
\end{tcolorbox}
